% Persistent Segment Tree

$O(\log N)$ time and space per update. \lstinline|upgrade(v, pos, val)| creates a new version from version \lstinline|v|. Query version \lstinline|v| via \lstinline|query(roots[v], 1, n, ql, qr)|.

\begin{lstlisting}
struct Node {
    int val; Node *l, *r;
    Node(int val = 0) : val(val), l(nullptr), r(nullptr) {}
    Node(Node* o) : val(o->val), l(o->l), r(o->r) {}
};

struct PersistentSegTree {
    int n; vector<Node*> roots;
    PersistentSegTree(int n) : n(n) { roots.push_back(build(1, n)); }

    Node* build(int l, int r) {
        Node* curr = new Node(0);
        if (l == r) return curr;
        int mid = (l + r) >> 1;
        curr->l = build(l, mid); curr->r = build(mid + 1, r);
        return curr;
    }

    Node* update(Node* prev_node, int l, int r, int pos, int val) {
        Node* curr = new Node(prev_node);
        if (l == r) { curr->val += val; return curr; }
        int mid = (l + r) >> 1;
        if (pos <= mid) curr->l = update(prev_node->l, l, mid, pos, val);
        else curr->r = update(prev_node->r, mid + 1, r, pos, val);
        curr->val = curr->l->val + curr->r->val;
        return curr;
    }

    int query(Node* node, int l, int r, int ql, int qr) {
        if (ql <= l && r <= qr) return node->val;
        int mid = (l + r) >> 1, res = 0;
        if (ql <= mid) res += query(node->l, l, mid, ql, qr);
        if (qr > mid) res += query(node->r, mid + 1, r, ql, qr);
        return res;
    }

    void upgrade(int version, int pos, int val) {
        roots.push_back(update(roots[version], 1, n, pos, val));
    }
};
\end{lstlisting}
